\part*{Part 0 --- The argument in ten pages}
\addcontentsline{toc}{part}{Part 0 --- The argument in ten pages}
\setcounter{section}{0}
\renewcommand{\thesection}{0.\arabic{section}}

\bigskip
\noindent
\emph{This part uses no category theory and states no theorems. It is the whole
argument, in the order in which it is best met for the first time.}
\bigskip

\section{The question}
\label{sec:0question}

Suppose you want to build something that learns about its world. Not a function
approximator fitted offline to a corpus, but a thing that is \emph{in} a world,
that spends something to look at it, that forms guesses and finds out whether
they were right, and that dies if it spends more than it gets back.

Very early in the design, before any question about architecture, a question
arrives that is easy to skip past: \emph{what can this thing say?} Whatever it
believes, it believes in some vocabulary. Whatever hypothesis it entertains, that
hypothesis is a sentence in some language. Whatever experiment it runs, the
experiment's outcome has to be expressible as a change in that language's truth
values. The vocabulary is not decoration on the learner. It is the shape of
everything the learner will ever be able to think.

The usual move is to pick one. Take first-order logic, or a probabilistic graphical
model, or a set of hand-designed features, and then get on with the interesting
work. The move is so standard that it barely registers as a move. But it has a
cost, and the cost is that the vocabulary is now unaccountable: it floats above the
machine, related to it by nothing in particular, and when it turns out to be the
wrong vocabulary there is no principled way to say what a better one would be.

The claim of this document is that the choice can be taken out of your hands, in a
good way. If you have already said what your computational system is --- what its
states look like and how it steps --- then you have \emph{already said} what its
logic is. The connectives are recoverable from the syntax. The modalities are
recoverable from the rewrite rules. There is nothing left to design. And what
falls out is not one logic but a structured range of them, because the recovery
procedure has a small number of parameters, and different settings of those
parameters give different --- genuinely, usefully different --- disciplines of
observation over the same underlying system.

We call these \emph{observation disciplines}, and the plan of this document is to
exhibit the range, show that some famous logics are points in it, and then argue
that one particular way of deploying such a discipline turns it from a description
into an instrument that a resource-bounded learner can actually hold.

\section{What a logic is made of}
\label{sec:0made}

Begin from a commitment that sounds technical but is really a discipline of
honesty: \emph{a formula denotes the collection of computations that witness it.}
This is realizability, in the sense that goes back to Kleene \cite{Kleene45}. It
says: do not tell me a formula is true; show me what makes it true, and let it be
a thing that could run.

Adopt that, and a logic falls into three parts. Not by stipulation --- by
inspection. Every connective you can build is of one of three kinds.

\paragraph{Structural connectives come from the term constructors.}
Your terms are built by constructors: a pairing, a send, a receive, a parallel
composition. Each of them lifts to a predicate about \emph{shape}. If terms can be
put side by side --- $P \mid Q$ --- then there is a connective $\varphi \mid \psi$
which holds of a term exactly when the term can be split into a part satisfying
$\varphi$ and a part satisfying $\psi$. If terms can send on a channel, there is a
connective saying \emph{is a send on a channel satisfying $\dots$}. These are not
invented; there is exactly one per constructor, and you get them all whether you
want them or not.

\paragraph{Behavioural connectives come from the rewrite rules.}
Your terms step. Each step happens somewhere: there is a surrounding context, and
a redex sitting in it. Label the step by that context and you have a modality:
$\langle K \rangle \varphi$ means \emph{there is a step, at the position marked by
$K$, after which $\varphi$ holds}. This is the Hennessy--Milner construction
\cite{HennessyMilner85} with contexts in place of action labels, which turns out
to matter a great deal: an action-labelled logic looking at a board game sees one
label, ``a move was made,'' while a context-labelled logic sees as many labels as
there are legal moves.

\paragraph{Collection connectives come from the container.}
A formula denotes a collection of witnesses --- but a collection \emph{of what
kind}? A set? A multiset? A list? A probability distribution? This is a real
choice, and it is the one that carries the propositional logic. If witnesses live
in sets, then two formulae can be intersected and united, and you get conjunction
and disjunction in their familiar form. If witnesses live in multisets, you can no
longer silently discard a duplicate, and what you get instead is a
\emph{multiplicative} conjunction that tracks how many times a resource was used.
The container's algebra is the logic's structural rules. This is not an analogy;
it is an identification, and Part I states it precisely.

\bigskip
\noindent
So a logic has three components, and two of them --- structural and behavioural
--- are read off the presentation of the system with no discretion at all. The
third is a genuine parameter. There is a fourth parameter, introduced below.

\section{Turning the dials}
\label{sec:0dials}

Here is the payoff of decomposing rather than choosing. The components are dials,
and turning them lands you on logics you already know, plus many you do not.

Turn the container dial from sets to multisets and the structural rules change:
you lose the right to use a witness twice, and what emerges is the multiplicative
fragment of \emph{linear logic}. Stack a powerset on top of a free commutative
monoid and add a distinguished element to serve as a pole, and what emerges is
classical linear logic in full, with its two conjunctions and its exponential
\cite{Girard87,Yetter90}. Part II works this instance carefully, because it is
the one that most convincingly makes the case: linear logic is famously the logic
of resources, and here it is not \emph{modelled} as a logic of resources but
\emph{generated} from a container that keeps count.

Set the container to the open sets of a topology and you get intuitionistic logic
in its geometric guise, where a proposition is an affirmation, something that can
be positively said \cite{Vickers89}. Now weaken the container by dropping one
axiom --- let the opens be closed under arbitrary unions but not under
intersection --- and you get Gabbay's \emph{semitopology} \cite{Gabbay25}, whose
points are participants and whose opens are coalitions with the resources to act
on their own. The intersection of two coalitions each strong enough to act need
not itself be strong enough to act. That single dropped axiom is the whole
difference from topology, and the logic it generates is the logic of decentralised
consensus. Conjunction does not survive as an operation; it survives as a
\emph{relation}, ``these two have a witness in common,'' and Gabbay arrived at
exactly that relation independently, for reasons of algebraic duality
\cite{GabbaySemiframes}. Two routes, one relation. That coincidence is the best
single piece of evidence that the generation story is tracking something real
rather than reorganising things we already knew how to write down.

The fourth dial is the one this document is most interested in, and it does not
appear in classical treatments at all: \emph{what does an observation return?}
The default answer is a bit --- true or false. But nothing forces it. Let a
formula return a number in $[0,1]$, or a non-negative real, or a value in a
tropical algebra where addition is minimum, or a complex amplitude. Each choice is
a coherent logic; each has different laws; and, as we will see, the choice
determines what a learner using that logic is capable of.

\begin{slogan}
A logic is a choice of observation discipline over a term language: what shapes
you may name, what steps you may see, how witnesses are collected, and what an
observation returns.
\end{slogan}

\begin{figure}[h]
\centering
\begin{tikzpicture}[font=\small,
  dial/.style={draw,rounded corners=2pt,align=center,inner sep=4pt,minimum width=3.15cm,minimum height=1.05cm},
  outbox/.style={draw,thick,rounded corners=2pt,align=center,inner sep=5pt,minimum width=3.2cm}]
\node[dial] (sig) at (0,1.35) {\textbf{signature} $\Sig$\\[1pt]
  \footnotesize what shapes you may name};
\node[dial] (rw) at (0,0) {\textbf{rewrites} $\Rw$\\[1pt]
  \footnotesize what steps you may see};
\node[dial] (co) at (0,-1.35) {\textbf{container} $\Coll$\\[1pt]
  \footnotesize how witnesses collect};
\node[dial] (v) at (0,-2.7) {\textbf{values} $\V$\\[1pt]
  \footnotesize what an observation returns};
\node[outbox] (L) at (5.6,-0.68) {an observation\\discipline\\[2pt]
  $\Logic(\Sig,\Rw,\Coll,\V)$};
\foreach \n in {sig,rw,co,v} { \draw[-{Stealth},gray] (\n.east) -- (L.west); }
\draw[dashed,gray] (-2.05,-0.675) -- (2.05,-0.675);
\node[anchor=east,font=\footnotesize,align=right] at (-2.2,0.15)
  {read off the\\presentation};
\node[anchor=east,font=\footnotesize,align=right] at (-2.2,-1.5)
  {genuine\\parameters};
\end{tikzpicture}
\caption{The four dials. The first two are read off the presentation of the
system; the last two are where the design freedom lives.}
\label{fig:dials}
\end{figure}

\section{Three places to put a formula}
\label{sec:0place}

Having a logic is not yet having an instrument. The same formula can be used in
three quite different ways, and only one of them makes it something a system can
look \emph{through}.

\paragraph{As a specification.} Write $\vdash \varphi \triangleright P$: from the
formula, synthesise a program that realises it. This is the proofs-as-programs
direction, and it is the direction most of the literature on logics for
concurrency takes. It is also the expensive direction, because it requires a proof
theory --- and for the logics generated here, the proof theory is usually the open
problem.

\paragraph{As a test.} Write $P \models \varphi$: given a program, check whether
the formula holds. This needs only the semantics, which the generation procedure
hands you for free. It is dramatically cheaper than synthesis, and it is what an
observer does.

\paragraph{As a guard.} Write the formula \emph{inside the program}, on the
binder:
\begin{center}
\texttt{for( y <- x where }$\varphi$\texttt{ )P}
\end{center}
Now the formula is evaluated at the moment a communication is about to happen,
against the candidate message, and its value decides whether the communication
occurs. This is the placement that makes a logic an instrument. The formula is no
longer a description of the computation; it is part of the computation, and it is
the part that decides what the computation does next.

That last point deserves emphasis, because it is the hinge of the whole document.
A concurrent system of any interest is non-deterministic: at any moment several
communications could occur, and something has to decide which does. In a
conventional implementation that decision is made by the hardware --- by which
core happened to win a race. It is made, in other words, by something the language
has nothing to say about. Putting a formula in the binder moves the decision
\emph{into the language}. The guard's value is what resolves the race. And since
the guard is a formula of the generated logic, the resolution of non-determinism
becomes a matter of \emph{what the system is able to observe about itself.}

\section{Why a Boolean guard cannot learn}
\label{sec:0boolean}

Now grade the guard, and watch what happens.

A crisp guard is a partition: it splits the candidate messages into those that
pass and those that do not. A graded guard is a map into a value algebra: each
candidate gets a number, and the numbers become propensities --- how likely, or
how strongly weighted, each possible communication is.

The difference sounds like a convenience. It is not. Consider a small foraging
organism from \cite{GradedWhere}. It lives in a world containing specimens of two
kinds: rich ones, worth opening, and decoys, which cost more to open than they
return. It has a cheap structural test $\chi$ --- a shape predicate, priced by its
depth --- which is \emph{evidence} rather than a decision procedure: rich specimens
usually pass it, decoys sometimes do. It has a limited budget. If it opens too many
decoys it dies.

Give it a crisp guard, and it opens whenever $\chi$ holds. In a generous world
(rich specimens common, the test informative) this works well: the crisp forager
achieves the highest terminal wealth of any strategy measured --- and the lowest
survival probability, $0.8501$. In a harsh world, where the base rate of rich
specimens drops to $0.15$ and the test gets noisier, the crisp forager's survival
probability is $0.0029$. It is, for practical purposes, extinct.

The graded forager survives at $0.9485$ in the first world and persists in the
second. The reason is not that it is cleverer. The reason is structural:

\begin{slogan}
A Boolean guard has no coordinate in which to record disappointment.
\end{slogan}

The crisp forager's guard is either satisfied or not. When it opens a decoy, that
outcome has nowhere to go. The formula $\chi$ was true; it is still true; nothing
about the guard has changed. To learn from the disappointment the crisp learner
would have to change the \emph{formula} --- move to a different, presumably deeper
and more expensive test --- which is a discrete jump, and an expensive one. The
graded forager's guard is the same formula, with a value attached, and the value
moves: over five openings, its belief slides from strength $0.750$ down to
$0.250$, and the probability that it opens the next specimen falls from $0.333$ to
$0.231$. It has learned to be careful without learning anything new.

This is the sense in which graded where-clauses are a learning instrument. The
formula language gives a learner a discrete space of hypotheses, and that space is
a tree: to move from one hypothesis to a nearby one you must back up and branch,
and there is a precise theorem \cite{ScientistNote} saying that a learner making
only small moves in that tree can never leave the region it started in. Grading
adds a continuous coordinate \emph{inside} each region. The picture that results
is not a space but a fibration: a tree of formulae, with a value-space hanging
over each point. Search in the tree is backtracking; search in the fibre is
something much closer to a gradient. A learner needs both, and until now the
framework had only one.

\section{Seeing is metered}
\label{sec:0metered}

The last ingredient is that observation costs something.

If the structural predicates are the learner's \emph{senses} --- and they are;
that is what it means for a formula to describe the shape of a term --- then using
them is work, and the work should be charged. The natural price is depth: a
predicate that looks one constructor deep costs less than one that looks five
deep, and a modality that requires actually driving the system a step costs more
than a predicate that merely reads structure already present. In the underlying
calculus this is not an add-on; there is a cost-accounting layer in which tokens
are conserved, computation is funded, and a computation that runs out of budget
stops \cite{ScientistNote,GsltOmnibus}.

Once observation is metered, sentences one could not previously state become
measurable. Here are three, all measured rather than argued.

\paragraph{Truth is not the best buy.} In a study of noughts-and-crosses players
whose strategies are rungs on a ladder of increasingly deep formulae
\cite{GameNote}, the player using the exactly-correct characteristic formula ---
the one that plays perfectly --- is \emph{dominated}. It wins $87.3\%$ of games at
a metered cost of $679$ units, while a shallower heuristic wins $91.8\%$ at $256$.
The correct theory is the only rung that never loses; it buys safety, not food.

\paragraph{Perception is cheaper than interaction, and it is where the yield is.}
On the same ladder, the three rungs that read structure the world has already
deposited --- with no modality at all --- deliver $92.3\%$ of the total gain for
$34.3\%$ of the total cost. The two rungs that require driving the system forward
deliver the remaining $7.7\%$ for the other $65.7\%$. There is a general thesis
here about environments rather than learners: a predicate can be \emph{read} if
the environment has already computed it into legible structure, or \emph{driven}
if it has not, and the interesting question is which party is billed.

\paragraph{Whether truth is affordable is a property of the world.} Run the same
ladder construction on a different game --- Nim \cite{ComposeNote} --- and the
result inverts. There the perfect strategy is the \emph{cheapest} rung on the
ladder and the best buy on it. Nothing about the learner changed. What changed was
the environment's structural regularity.

\section{The range, in one page}
\label{sec:0range}

Six disciplines are worked in Part VI. Each is a different setting of the dials,
and each comes with a measured example.

\begin{enumerate}[leftmargin=2em]
\item \textbf{Spatial, crisp.} Only structural predicates, Boolean values. The
senses. Worked on board games: forks, threats, and the discovery that requiring a
fork to be expressible forces the encoding of the board.

\item \textbf{Modal, crisp.} Add the context-labelled modality: lookahead. Worked
on the same games, where it turns out to be the most expensive and least
productive part of the ladder.

\item \textbf{Graded, real-valued.} Values in $[0,1]$ or $\R_{\ge 0}$, supplied by
a probabilistic logic. The mortal forager, and the discovery that the personality
parameter governing scepticism has an interior metabolic optimum --- credulous
foragers open decoys and die, sceptical ones starve, and the optimum rises
roughly sevenfold with the risk of the world.

\item \textbf{Graded, complex-valued.} Amplitudes rather than probabilities.
Interference becomes available, which means Shor's algorithm can be written in the
guard language --- and, more informatively, the semantics turns out to be
\emph{too} powerful rather than too weak, so that a discipline restricting clauses
to be norm-preserving is what makes it a quantum language rather than a magic one.

\item \textbf{Semitopological.} The container is the opens of a semitopology, and
the logic is about coalitions. Worked twice: on a hunting pride reading which way
a herd will break, and on a honeybee swarm choosing a nest site. From the swarm
comes the cleanest prediction in the corpus --- that two rival quorums require
twice the quorum threshold in scouts, so that a swarm smaller than that cannot
split no matter what, whatever the site qualities or the inhibition dynamics.

\item \textbf{Resolution-graded.} The observation discipline is coarsened
deliberately, and one asks what an observer at that coarseness reports. It
reports fewer kinds than the world contains, and by a computable factor: at the
coarsest resolution, one kind where there are four, and a population count
inflated by $3.4\times$.
\end{enumerate}

That last item is where the document ends, because it is where the argument
becomes epistemology rather than engineering. Choosing an observation discipline
is choosing which distinctions exist \emph{for you}. A learner does not fail to
see the distinctions it cannot afford; it fails to see that they are there, and
reports a world with fewer kinds in it than the world has. That is a Kantian
thought with a ledger attached, and the ledger is the contribution.

\section{The evidence, on one page}
\label{sec:0evidence}

Almost everything asserted above has been measured. Table~\ref{tab:collected0}
collects the headline figures; Part VI works each of them, and the simulations
that produced them ship with the notes cited in the last column. A reader who
stops at the end of Part 0 should at least have seen this table.

\begin{table}[h]
\centering
\small
\renewcommand{\arraystretch}{1.25}
\begin{tabular}{@{}>{\raggedright\arraybackslash}p{5.4cm}>{\raggedright\arraybackslash}p{4.4cm}>{\raggedright\arraybackslash}p{4.2cm}@{}}
\toprule
result & figure & source \\
\midrule
Non-modal rungs dominate the ladder
  & $92.3\%$ of yield for $34.3\%$ of cost
  & \cite{GameNote} \\
The correct theory is dominated
  & $.873$ at $679.2$ vs $.918$ at $255.9$
  & \cite{GameNote} \\
Confinement, measured
  & best opening reverses; $.833 \to .548$
  & \cite{GameNote} \\
Affordability inverts with the world
  & Nim returns rising: $+.00602/+.01262/+.01476$
  & \cite{ComposeNote} \\
Bootstrap trap
  & $0.00$ openings, $520.0$ vs $3543.7$
  & \cite{GradedWhere} \\
Crisp guard cannot learn to stop
  & survival $.0029$ in the harsh world
  & \cite{GradedWhere} \\
Confidence buys variance, not yield
  & $.8501/4084.7$ vs $.9485/3543.7$
  & \cite{GradedWhere} \\
Scepticism has an interior optimum
  & $k^{*} = 9$; ensemble $k^{*} = 64$
  & \cite{GradedWhere} \\
Interleaving inflation is real
  & exactly $2$, $6$, $24$ at $n=2,3,4$
  & \cite{GradedWhere} \\
Shor as deliberate interference deficit
  & $12$ of $16$ outcomes at amplitude $0$
  & \cite{GradedWhere} \\
The modality buys irreversibility
  & precision $1.0000$ at $0.37\%$ coverage
  & \cite{QuorumNote} \\
Illegibility denies liveness, not safety
  & coverage $.7520 \to .0115$, precision held
  & \cite{QuorumNote} \\
The margin is a clock
  & $42.4/29.6/28.9/18.0$ ticks by quartile
  & \cite{QuorumNote} \\
The $2k$ theorem
  & split $0.0000$ below $S=30$; $0.0130$ at $30$
  & \cite{QuorumNote} \\
Cross-inhibition changes jobs
  & deadlock $.5120\to.0000$; split $.3860\to.1170$
  & \cite{QuorumNote} \\
A coarse observer under-reports kinds
  & $1$ kind of $4$; count inflated $3.4\times$
  & \cite{ScopeNote} \\
\bottomrule
\end{tabular}

\caption{The headline measurements. Worked in Part VI.}
\label{tab:collected0}
\end{table}

\section{What is claimed, and what is not}
\label{sec:0claims}

\paragraph{Claimed.} That the structural and behavioural components of a logic
over a term language are determined by the presentation of that language, with no
design freedom; that the remaining freedom is concentrated in two parameters, the
collection container and the value algebra; that several established logics ---
intuitionistic, classical linear, coalitional --- are settings of those
parameters; that placing such a formula in a binder's guard makes it the
resolution mechanism for the system's own non-determinism; that grading the guard
gives a resource-bounded learner a continuous coordinate it otherwise lacks; and
that all of this is priced, so that the resulting claims about what a learner can
afford to see are measurable and have been measured.

\paragraph{Not claimed.} That every generated logic is well-behaved: adequacy ---
the property that formulae separate exactly what interaction separates ---
depends on the target logic having enough structure, and many useful settings do
not have it. That the generated logic is decidable: it generally is not, and the
practical discipline is to restrict either the language fragment or the connective
set until model-checking terminates, at the cost of adequacy. That the results
about foragers, prides and swarms are established biology: the swarm predictions
in particular are untested proposals, offered because they are falsifiable and
cheap to test, not because they have been confirmed. That this document proves
much: it is a summary, and where it advances a new claim it is normally left as a
conjecture, marked as such.

\paragraph{One methodological note.} Almost every number quoted in this document
comes from a simulation written alongside the note it appears in, and those
simulations are in the same repository as the notes. Where a result is measured
rather than proved, it is said to be measured. Where a result is a negative
result --- something the framework predicted that did not survive contact with the
arithmetic --- it is kept. There are several.
